\section*{अध्याय \ref{ch:b1:populations} --- आबादियाँ और जनांकिकी}

\begin{solution}{exo:b1:populations:1}
\omterm{def:b1:populations:population}{आबादी}: एक ही जगह और समय में एक ही जाति के प्राणी,
जो आपस में प्रजनन कर सकें। घनत्व: प्रति इकाई क्षेत्रफल या आयतन प्राणी।
परिक्षेपण: उनका दिक्-ढर्रा --- गुच्छित (कोई झुंड, बिच्छू-बूटी का कोई
धब्बा), एकसमान (एक चोंच की दूरी पर घोंसले बनाते गैनेट, क्रिओसोट झाड़ियाँ),
यादृच्छ (किसी लॉन पर सिंहपर्णी)।
\end{solution}

\begin{solution}{exo:b1:populations:2}
$l_x$: आयु $x$ पर जीवित समवर्ग का अंश। $m_x$: आयु $x$ की प्रति
मादा बेटियाँ। $R_0 = \sum l_x m_x$: किसी नवजात मादा की जीवन भर की
बेटियाँ। $T = \sum x l_x m_x / R_0$: माताओं की माध्य आयु।
$R_0 = 1$: हर मादा अपनी जगह भर देती है; \omterm{def:b1:populations:population}{आबादी} स्थिर है।
\end{solution}

\begin{solution}{exo:b1:populations:3}
प्रकार III: 1000 में से लगभग 15, यानी \qty{1.5}{\%}। प्रकार I: 1000 में
से लगभग 700, यानी \qty{70}{\%}।
\end{solution}

\begin{solution}{exo:b1:populations:4}
$\dd N/\dd t = rN$ और $\dd N/\dd t = rN(1 - N/K)$। $r$ वह
प्रति-व्यक्ति वृद्धि दर है जब संसाधन असीम हों (प्रति प्राणी प्रति इकाई
समय जन्म घटा मृत्यु); और $K$ \omterm{thm:b1:populations:logistic}{वहन क्षमता} है, यानी
वह आकार जिस पर वृद्धि रुक जाती है।
\end{solution}

\begin{solution}{exo:b1:populations:5}
$r = \ln(1000)/5 = \qty{1.38}{h^{-1}}$; द्विगुणन $\ln 2/1.38 =
\qty{0.5}{h}$।
\end{solution}

\begin{solution}{exo:b1:populations:6}
$l_x m_x = 0, 0.5, 0.6, 0.2, 0$: $R_0 = 1.3$। $T = (1\times 0.5 +
2\times 0.6 + 3\times 0.2)/1.3 = 2.3/1.3 = 1.77$। $r = \ln 1.3/1.77 =
\qty{0.15}{}$ प्रति इकाई समय: वह बढ़ती है।
\end{solution}

\begin{solution}{exo:b1:populations:7}
$N = 50\times 70/7 = 500$। केवल 40 चिह्नित मछलियाँ बचीं: $N = 40\times
70/7 = 400$; पहला आकलन एक-चौथाई से बहुत ऊँचा था (\omterm{def:b1:populations:population}{आबादी} का चिह्नित
अंश माने गए से छोटा था)।
\end{solution}

\begin{solution}{exo:b1:populations:8}
$0.2\times 100\times 0.9 = 18$; $0.2\times 500\times 0.5 = 50$;
प्रति वर्ष $0.2\times 900\times 0.1 = 18$। अधिकतम $rK/4 = 50$ प्रति वर्ष,
$N = 500$ पर।
\end{solution}

\begin{solution}{exo:b1:populations:9}
सिंहपर्णी: $r$ (हज़ारों बीज, वार्षिक)। हाथी: $K$ (चार साल में एक
बछड़ा, लंबा जीवन, देखभाल)। कॉड: $r$ (लाखों अंडे, कोई देखभाल नहीं)।
गोरिल्ला: $K$ (चार साल में एक शिशु, दशकों का जीवन)। घरेलू मक्खी: $r$
(सैकड़ों अंडे, दस दिन में परिपक्व)।
\end{solution}

\begin{solution}{exo:b1:populations:10}
$r = \ln(29/10)/2 = \qty{0.53}{h^{-1}}$; $K \approx 660$। \qty{8}{h} पर:
$660/(1 + 65e^{-4.24}) = 660/(1 + 0.94) = 340$, जबकि गिने गए 351.
सबसे तेज़ $K/2 = 330$ पर, जो $t = \ln 65/0.53 = \qty{7.9}{h}$ पर आता है,
जहाँ $\dd N/\dd t = rK/4 = 87\times 10^5$ \omterm{def:b1:cell-unit-of-life:cell}{कोशिकाएँ} प्रति मिलीलीटर प्रति
घंटा।
\end{solution}

\begin{solution}{exo:b1:populations:11}
नियमन के लिए ज़रूरी है कि $N$ चढ़ने पर प्रति-व्यक्ति दर गिरे, ताकि
किसी विक्षोभ के बाद $N$ किसी साम्य की ओर लौटे; जो कारक
एक नियत अंश मारता है वह $N$ बदलता है पर $N$ पर दर की निर्भरता
नहीं, इसलिए उसके बाद \omterm{def:b1:populations:population}{आबादी} वही वृद्धि फिर शुरू कर देती है।
मौसम कीटों को इस तरह सीमित करता है कि वह इतनी बार और इतने यादृच्छ ढंग से
चोट करता है कि \omterm{def:b1:populations:population}{आबादी} किसी छत के पास पहुँचने से पहले ही पीछे धकेल दी जाती
है: यानी बार-बार के विक्षोभ से दबी रखी कोई अनियमित \omterm{def:b1:populations:population}{आबादी}, जिसका
औसत आकार किसी \omterm{thm:b1:populations:logistic}{वहन क्षमता} के बजाय बुरे मौसमों की बारंबारता पर निर्भर करता
है।
\end{solution}

\begin{solution}{exo:b1:populations:12}
एक ही आकार की दो \omterm{def:b1:populations:population}{आबादियाँ}, जिनमें से एक का युवा आधार चौड़ा हो और दूसरी के
आयु-वर्ग बराबर, अलग-अलग भविष्य रखती हैं: पहली दशकों तक बढ़ती रहेगी
भले ही उसका $m_x$ अभी प्रतिस्थापन पर गिर जाए, क्योंकि उसके
बहुत-से युवाओं को अभी जनन करना है (जनांकिकीय संवेग); दूसरी
नहीं बढ़ेगी। $N$ अतीत दर्ज करता है; आयु-बंटन और उस पर
$m_x$ की अनुसूची ही वे हैं जिन पर समीकरण काम करते हैं, और
पिरामिड अगली पीढ़ी का आकार बताता है जबकि $N$ उसके बारे में कुछ नहीं
कहता।
\end{solution}

\begin{solution}{pb:b1:populations:1}
\textbf{1.} $r = \ln(29/10)/2 = \qty{0.53}{h^{-1}}$।
\textbf{2.} $\ln 2/0.53 = \qty{1.3}{h}$।
\textbf{3.} $K \approx 660\times 10^5 = \qty{6.6e7}{}$ \omterm{def:b1:cell-unit-of-life:cell}{कोशिकाएँ} प्रति
मिलीलीटर।
\textbf{4.} $N(t) = 660/(1 + 65e^{-0.53t})$: 4 h पर, $660/(1 + 7.8) =
75$ (गिने 71); 8 h पर, 340 (351); 12 h पर, $660/(1 + 0.112) = 594$
(595)।
\textbf{5.} $K/2 = 330$ पर, जो $t = \ln 65/0.53 = \qty{7.9}{h}$ पर आता है;
अधिकतम दर $rK/4 = 0.53\times 660/4 = 87\times 10^5$ \omterm{def:b1:cell-unit-of-life:cell}{कोशिकाएँ} प्रति
मिलीलीटर प्रति घंटा।
\textbf{6.} $10e^{0.53\times 12} = 10\times 578 = 5780$: यानी देखे गए
595 का दस गुना।
\textbf{7.} शर्करा (या ऑक्सीजन) का चुक जाना; और इथेनॉल तथा अम्लों का
जमाव जो \omterm{def:b1:cell-unit-of-life:cell}{कोशिकाओं} को विषाक्त करता है।
\textbf{8.} $q_0 = 0.40$; $q_4 = 1 - 0.34/0.40 = 0.15$: यानी भारी शुरुआती
मर्त्यता फिर एक ठहरी दर --- प्रकार II और III के बीच, जैसा जंगल में
अधिकांश बड़े स्तनधारियों के लिए होता है।
\textbf{9.} $l_x m_x = 0, 0, 0.312, 0.368, 0.320, 0.272, 0.182, 0.080,
0.018, 0$: $R_0 = 1.55$।
\textbf{10.} $\sum x l_x m_x = 0.624 + 1.104 + 1.280 + 1.360 + 1.092 +
0.560 + 0.144 = 6.16$; $T = 6.16/1.55 = \qty{4.0}{yr}$।
\textbf{11.} $r = \ln 1.55/4.0 = \qty{0.11}{yr^{-1}}$; द्विगुणन
\qty{6.3}{yr} में।
\textbf{12.} $200e^{1.1} = 600$ हिरन।
\textbf{13.} $N(10) = 600/(1 + 2e^{-1.1}) = 600/1.67 = 360$; और 500 तब
पार होता है जब $(K - N_0)/N_0\,e^{-rt} = 0.2$ हो, यानी $e^{-rt} = 0.1$, $t = \ln
10/0.11 = 21$ वर्ष।
\textbf{14.} $120\times 150/18 = 1000$ पर्च।
\textbf{15.} 100 चिह्नित: $100\times 150/18 = 833$; पहला आकलन
बहुत ऊँचा था (माने गए से कम चिह्नित मछलियाँ बाहर थीं)।
\textbf{16.} $rK/4 = 100$ प्रति वर्ष, $N = 500$ पर।
\textbf{17.} $T_2 = \ln 2/0.4 = \qty{1.7}{yr}$।
\textbf{18.} पकड़ $hN$, जिसमें $h = 0.2$: $rN(1 - N/K) = hN$ से $N^* =
K(1 - h/r) = 1000\times 0.5 = 500$ मिलता है और प्रति वर्ष \num{100} की
पकड़ --- यानी अधिकतम सतत उपज से नीचे, पर कोई अंश अपने आप भंडार के हिसाब
से बैठ जाता है: यदि भंडार आधा हो जाए तो पकड़ भी आधी हो जाती है, और \omterm{def:b1:populations:population}{आबादी}
लौट आती है।
\textbf{19.} 500 पर: $0.4\times 500\times 0.5 = 100 < 120$: \omterm{def:b1:populations:population}{आबादी}
घटती है। 300 पर: $0.4\times 300\times 0.7 = 84 < 120$: और तेज़ी से
घटती है --- $K/2$ से नीचे आते ही अधिकतम उपज से ऊपर की कोई नियत कटाई
उसे विलुप्ति तक ले जाती है।
\textbf{20.} $K/2$ पर अधिशेष अपने अधिकतम पर है पर कोई भी त्रुटि ---
पाँचवें भाग से ऊँची आँकी गई कोई \omterm{def:b1:populations:population}{आबादी}, कोई बुरा साल जो $r$ काट दे ---
किसी सतत पकड़ को गिरावट में बदल देती है, और $K/2$ से नीचे \omterm{def:b1:populations:population}{आबादी} के साथ
अधिशेष भी सिकुड़ता है: यानी वह साम्य अति-कटाई के प्रति अस्थायी है।
$K/2$ से कुछ ऊपर कटाई करने से एक अंतर बचा रहता है।
\textbf{21.} $R_0$ आधा होकर $0.78 < 1$ हो जाता है: झुंड तब तक सिकुड़ता है
जब तक वह 400 से नीचे न आ जाए और $m_x$ उबर न आए --- यानी कोई
घनत्व-निर्भर कारक, जो झुंड को 400 के पास नियमित करता है।
\textbf{22.} नहीं: वह आकार चाहे जो हो \qty{40}{\%} हटा देता है और
प्रति-व्यक्ति दर अपरिवर्तित छोड़ देता है; आगे के वर्षों में झुंड
किसी नीची शुरुआत से उसी $r$ पर फिर बढ़ आता है --- यानी कोई विक्षोभ, कोई
नियमन नहीं।
\textbf{23.} $0.11N(1 - N/600) = 30$: $N^2 - 600N + 163\,600 = 0$,
$N = 300\pm\sqrt{90\,000 - 163\,600}$ --- कोई वास्तविक हल नहीं: प्रति वर्ष
30 अधिकतम अधिशेष $rK/4 = 16.5$ से अधिक है, और भेड़िए
झुंड को विलुप्ति तक ले जाते हैं। प्रति वर्ष 15 के साथ: $N = 300\pm\sqrt{90\,000 -
81\,800} = 300\pm 91$, यानी 209 और 391; ऊपर वाला स्थायी है
(उससे ऊपर अधिशेष 15 से कम है और झुंड वापस गिरता है, उससे नीचे अधिक और
वह चढ़ता है), और नीचे वाला अस्थायी।
\textbf{24.} लिंक्स की संख्या भोजन पर विलंब से जवाब देती है: खरगोशों की
बढ़त अगले एक-दो साल भर लिंक्स के जन्म तथा \omterm{def:b1:populations:lifetable}{उत्तरजीविता} चढ़ाती है, और
खरगोशों का ढहना लिंक्स को तभी भूखा मारता है जब उनकी चरबी और उस मौसम के
बच्चे चुक जाएँ।
\textbf{25.} यीस्ट: $K = \qty{6.6e7}{}$ \omterm{def:b1:cell-unit-of-life:cell}{कोशिकाएँ} प्रति मिलीलीटर, $r =
\qty{0.53}{h^{-1}}$। हिरन: $R_0 = 1.55$, $T = \qty{4.0}{yr}$, $r =
\qty{0.11}{yr^{-1}}$।
\end{solution}
